\def\level{Première}  % Classe à droite
\def\course{NSI}
\def\chaptername{Constructions élémentaires} % Titre au centre
\def\docname{LP01}        % Identifiant à gauche

\pagestyle{cours_FP}
\makeheader{} % Affiche le bandeau


\section{Les variables en Python}

\begin{easybox}{Déclarer une variable en Python}

	Pour déclarer une variable en Python, il suffit de lui donner un nom et de lui affecter une valeur à l'aide du symbole $=$. Par exemple, pour déclarer une variable \texttt{x} qui contient la valeur 10, on écrit:

\begin{minipage}{0.48\linewidth}
	\begin{CodePythontexAlt}[Largeur=1\linewidth,Centre,Lignes=true,PremLigne=1, Lignes=false,]{}
    a = 10
    \end{CodePythontexAlt}
\end{minipage}\hfill{}
\begin{minipage}{0.48\linewidth}
\begin{scratch}
\blocklook{Affecter à \selectmenu{x} la valeur \ovalnum{10}}\end{scratch}
\end{minipage}
\end{easybox}


\begin{easybox}{Le nom des variables en Python}
	\begin{itemize}
		\item Elles doivent commencer par une lettre ou un underscore (\texttt{\_}).
		\item Elles peuvent contenir des lettres, des chiffres et des underscores.
		\item Elles ne peuvent pas être des mots réservés du langage Python (comme \texttt{if}, \texttt{for}, \texttt{while}, etc.).
		\item \texttt{Variable} et \texttt{variable} sont considérées comme deux variables différentes.
	\end{itemize}
\end{easybox}

\begin{exemple}[Noms de variables valides et invalides]{}

	  \renewcommand{\arraystretch}{1.}
    \begin{tabular}{|l|p{12cm}|}
    \hline
    Bonne variable & \verb+ma_variable+ \\ \hline
    Mauvaise variable & \verb+Ma_variable+,  \verb+ma variable+, \verb+1_variable+, \verb+ variable@+ \\ \hline
    A éviter    & \verb+Variable non explicites+ \\ \hline
    \end{tabular}
\end{exemple}

\begin{easybox}{Afficher une variable en Python}

	Pour afficher la valeur d'une variable en Python, on utilise la fonction \texttt{print()}. Par exemple, pour afficher la valeur de la variable \texttt{x}, on écrit:

\begin{minipage}{0.48\linewidth}
	\begin{CodePythontexAlt}[Largeur=1\linewidth,Centre,Lignes=true,PremLigne=1, Lignes=false,]{}
    print(x)
    \end{CodePythontexAlt}
\end{minipage}\hfill{}
\begin{minipage}{0.48\linewidth}
\begin{scratch}
\blocklook{Afficher la valeur de \selectmenu{x}}\end{scratch}
\end{minipage}
\end{easybox}

\begin{easybox}{Les types de variables en Python}
	\begin{itemize}
		\item \textbf{int}: pour les nombres entiers (ex: 1, 42, -5)
		\item \textbf{float}: pour les nombres à virgule flottante (ex: 3.14, -0.001)
		\item \textbf{str}: pour les chaînes de caractères (ex: "Bonjour", 'Python')
		\item \textbf{bool}: pour les valeurs booléennes (ex: True, False)
	\end{itemize}
\end{easybox}

\begin{minipage}{0.19\linewidth}
\begin{lstlisting}[numbers=none]
# type (int)
a = 5
b = 6
\end{lstlisting}
    \end{minipage}\hfill
    \begin{minipage}{0.19\linewidth}    
\begin{lstlisting}[numbers=none]   
# type (str)
a = 'Bonjour'
b = '6'   
\end{lstlisting}    
    \end{minipage}\hfill
    \begin{minipage}{0.19\linewidth}    
\begin{lstlisting}[numbers=none] 
# type (float)
a = 5.0 
b = 6.8
\end{lstlisting}    
    \end{minipage}\hfill{}
\begin{minipage}{0.19\linewidth}    
\begin{lstlisting}[numbers=none]  
# type (bool)
a = True
b = False
\end{lstlisting}    
    \end{minipage}

\begin{easybox}{Opérations mathématiques de base en Python}{}
	\begin{minipage}{0.43\linewidth}
		\begin{itemize}[]
        \item[] \texttt{+} \quad  addition:  \verb~5 + 3~ donne 8
        \item[] \texttt{-} \quad  soustraction \verb~5 - 3~ donne 2
        \item[] \texttt{*} \quad  multiplication \verb~5 * 3~ donne 15
        \item[] \texttt{/} \quad  division \verb~6 / 2~ donne 3.0 (\texttt{float})
\end{itemize}
	\end{minipage}\hfill{}
	\begin{minipage}{0.55\linewidth}
\begin{itemize}[]
        \item[] \texttt{**} \quad  puissance \verb~2 ** 3~ donne 8
        \item[] \texttt{\%} \quad  reste de la division entière \texttt{7 \% 3} donne 1
        \item[] \texttt{//} \quad  quotient de la division entière \texttt{7 // 3} donne 2
    \end{itemize}
	\end{minipage}
\end{easybox}

\begin{minipage}{0.3\linewidth}
    \begin{ConsolePythontex}[Label=false]{}
        print(17+2)
        print(17-2)
    \end{ConsolePythontex}
\end{minipage}\hfill
\begin{minipage}{0.3\linewidth}
    \begin{ConsolePythontex}[Label=false]{}
		print(17*2)
		print(17/2)
    \end{ConsolePythontex}
\end{minipage}\hfill
\begin{minipage}{0.3\linewidth}
    \begin{ConsolePythontex}[Label=false]{}
        print(17%2)
        print(17//2)
    \end{ConsolePythontex}
\end{minipage}

\begin{exercice_cours}[Repérer les les types de variables et leurs valeurs dans les programmes suivants]
	\begin{minipage}{0.48\linewidth}
		\begin{itemize}
		\item \verb~x = 5~
		\item \verb~y = "Hello"~
		\item \verb~z = True~
		\item \verb~a = 3.14~
	\end{itemize}
	\end{minipage}\hfill{}
	\begin{minipage}{0.48\linewidth}
	\begin{itemize}
		\item \verb~c = 6/2~
		\item \verb~d = 6//2~
		\item \verb~f = '1' + '2'~
		\item \verb~g = 2**3~
	\end{itemize}
\end{minipage}
	
\end{exercice_cours}


\medskip

\section{Les variables de type chaine de caractère \texttt{str}}


\begin{easybox}{Les chaines de caractères}{}
    \smallskip
    Les chaines de caractères sont des variables qui contiennent du texte. Elles sont délimitées par des guillemets simples ou doubles.

    \begin{minipage}{0.32\linewidth}
        \begin{ConsolePythontex}[ Label=false]{}
            a = 'Bonjour'
            b = ' a tous'
            c = a + b
            print(c)
        \end{ConsolePythontex}
    \end{minipage}\hfill
    \begin{minipage}{0.32\linewidth}
        \begin{ConsolePythontex}[ Label=false]{}
            a = 'Bonjour'
            b = ' a tous'
            c = a * 3
            print(c)
        \end{ConsolePythontex}
    \end{minipage}\hfill
    \begin{minipage}{0.32\linewidth}
        \begin{ConsolePythontex}[ Label=false]{}
            a = 'Bonjour'
            b = ' a tous'
            c = a[0]
            print(c)
        \end{ConsolePythontex}
    \end{minipage}

    \medskip

    Dans l'exemple ci-dessus, la variable \verb+c+ contient la concaténation de \verb+a+ et \verb+b+ dans le premier cas, la répétition de \verb+a+ trois fois dans le deuxième cas et le premier caractère de \verb+a+ dans le troisième cas.
\end{easybox}

\begin{easybox}{indexation des chaines de caractères}{}
    
    \begin{minipage}{0.5\linewidth}
 Les chaines de caractères peuvent être vues comme un ensemble de caractères ``collés`` les uns aux autres

    Chaque caractère d'une chaine de caractères est indexé: 
    il y a un index qui donne sa position dans la chaine       
    \end{minipage}\hfill{}
    \begin{minipage}{0.5\linewidth}
 \begin{center}
        \begin{tabular}{|c|c|c|c|c|c|c|c|}
        \hline
        \textbf{Indice} & 0 & 1 & 2 & 3 & 4 & 5 & 6 \\
        \hline
        \textbf{Caractère} & B & o & n & j & o & u & r \\
        \hline
        \textbf{Indice Négatif} & -7 & -6 & -5 & -4 & -3 & -2 & -1 \\
        \hline
        \end{tabular}
    \end{center}       
    \end{minipage}
    
   
    

    \begin{minipage}[t]{0.48\linewidth}
        \begin{ConsolePythontex}[ Label=false]{}
            mot = 'Bonjour'
            mot[0]
            mot[4]
        \end{ConsolePythontex}
    \end{minipage}\hfill
    \begin{minipage}[t]{0.48\linewidth}
        \begin{ConsolePythontex}[ Label=false]{}
mot = 'Bonjour'
n = len(mot)
n
mot[-1]
		\end{ConsolePythontex}
    \end{minipage}
   
\end{easybox}

\begin{easybox}{Les opérations sur les chaines de caractères}{}   
    \smallskip
    \begin{itemize}
        \item[] \texttt{+} \quad pour la concaténation:  \verb~'Bonjour' + ' a tous'~ donne \verb~'Bonjour a tous'~
        \item[] \texttt{*} \quad pour la répétition:  \verb~'Bonjour' * 3~ donne \verb~'BonjourBonjourBonjour'~
        \item[] \texttt{[]} \quad pour l'indexation:  \verb~'Bonjour'[0]~ donne \verb~'B'~
        \item[] \texttt{in} \quad pour le test d'appartenance:  \verb~'on' in 'Bonjour'~ donne \verb~True~
        \item[] \texttt{not in} \quad pour le test d'appartenance:  \verb~'on' not in 'Bonjour'~ donne \verb~False~
        \item[] \texttt{len()} \quad pour la longueur:  \verb~len('Bonjour')~ donne \verb~7~
       
    \end{itemize}
\end{easybox}


\begin{minipage}[t]{0.32\linewidth}
    \begin{ConsolePythontex}[ Label=false]{}
        mot1 = 'Bonjour'
        mot2 = 'Thomas'
        mot1 + ' ' + mot2
    \end{ConsolePythontex}
\end{minipage}\hfill
\begin{minipage}[t]{0.32\linewidth}
    \begin{ConsolePythontex}[ Label=false]{}
        'j' in 'Bonjour'
        'b' in 'Bonjour'
    \end{ConsolePythontex}
\end{minipage}\hfill{}
\begin{minipage}[t]{0.32\linewidth}
    \begin{ConsolePythontex}[ Label=false]{}
        'abracadabra'[4]
        len('abracadabra')
    \end{ConsolePythontex}
\end{minipage}

\begin{easybox}{Opérations avancées sur les chaines de caractères: le Slicing}{}
\bigskip
    \renewcommand{\arraystretch}{1.2}
    \begin{tabular}{p{2cm}|p{13cm}}
        \texttt{[a:b]} & conserve les caractères d'indice \texttt{[a:b]}:  \verb~'Bonjour'[1:4]~ donne \verb~'onj'~ \\
        \texttt{[a:]} & conserve les caractères d'indice \texttt{[a:]}:  \verb~'Bonjour'[3:]~ donne \verb~'jour'~ \\
        \texttt{[:b]} & conserve les caractères d'indice \texttt{[:b]}:  \verb~'Bonjour'[:3]~ donne \verb~'Bon'~ \\
        \texttt{[::-1]} & inverse les caractères:  \verb~'Bonjour'[::-1]~ donne \verb~'ruojnoB'~ \\
        \texttt{[1:-1]} & élimine le premier et dernier caractère:  \verb~'Bonjour'[1:-1]~ donne \verb~'onjou'~ \\
    \end{tabular}
\end{easybox}

\begin{exercice_cours}[Prédire ce que vont afficher les codes suivants]{}

    \begin{minipage}[t]{0.32\linewidth}        
\begin{lstlisting}[numbers=none]
>>> mot = 'Numerique'
>>> n = len(mot)
>>> mot[5]
>>> mot[n]
>>> mot[-1]
>>> mot[8]
\end{lstlisting}
    \end{minipage} \hfill
    \begin{minipage}[t]{0.32\linewidth}        
\begin{lstlisting}[numbers=none]
>>> mot = 'Numerique'
>>> m in mot
>>> e not in mot
>>> mot+mot[2]
>>> mot*2
>>> 'abc'+mot
\end{lstlisting}
    \end{minipage}\hfill
\begin{minipage}[t]{0.32\linewidth}        
\begin{lstlisting}[numbers=none] 
>>> mot = 'Numerique'
>>> mot[1:5]
>>> mot[:4]
>>> mot[3:]
>>> mot[1:-1]
\end{lstlisting}
    \end{minipage}
\end{exercice_cours}

\begin{easybox}{Principe du transtypage}{}
    Les types des variables en python sont dynamiques, on peut modifier sa valeur ainsi que son type.
    
    \begin{minipage}[t]{0.48\linewidth}
        \begin{ConsolePythontex}[ Label=false]{}
            a = '5'     # a est de type str
            b = int(a)  # b est de type int
            print(b)
        \end{ConsolePythontex}
    \end{minipage}\hfill
    \begin{minipage}[t]{0.48\linewidth}
        \begin{ConsolePythontex}[ Label=false]{}
            a = 4       # a est de type int
            b = str(a)  # b est de type str
            print(b)
        \end{ConsolePythontex}
    \end{minipage}
    
    \end{easybox}

    \section{Les variables de type booléen \texttt{bool}}

    
\begin{easybox}{Les booléens}{}
    \smallskip
    Les booléens sont des variables qui ne peuvent prendre que deux valeurs: \verb+True+ (vrai) ou \verb+False+ (faux).
    
    \smallskip
    Les booléens peuvent être utilisés dans des opérations logiques.
    
	\medskip

    \begin{minipage}{0.32\linewidth}
        \begin{ConsolePythontex}[ Label=false]{}
a = True
b = False
c = a and b
c
        \end{ConsolePythontex}
    \end{minipage}\hfill
    \begin{minipage}{0.32\linewidth}
        \begin{ConsolePythontex}[ Label=false]{}
a = True
b = False
c = a or b
c
        \end{ConsolePythontex}
    \end{minipage}\hfill
    \begin{minipage}{0.32\linewidth}
        \begin{ConsolePythontex}[ Label=false]{}
a = True
c = not a
c
        \end{ConsolePythontex}
    \end{minipage}    
\end{easybox}

\medskip

\begin{minipage}[t]{6cm}
    \begin{center}Table de vérité du \verb+not+\end{center}
    \begin{longtable}[]{@{}cc@{}}
        \hline Expression & Résultat\tabularnewline\hline\endhead{}
        \texttt{not} True & False\tabularnewline{}
        \texttt{not} False & True\tabularnewline\hline{}
    \end{longtable}
\end{minipage}
\begin{minipage}[t]{6cm}
    \begin{center}Table de vérité du \verb+and+\end{center}
    \begin{longtable}[]{@{}cc@{}}
        \hline Expression & Résultat\tabularnewline\hline\endhead{}
        True \texttt{and} True & True\tabularnewline{}
        True \texttt{and} False & False\tabularnewline{}
        False \texttt{and} True & False\tabularnewline{}
        False \texttt{and} False & False\tabularnewline\hline{}
    \end{longtable}
\end{minipage}
\begin{minipage}[t]{6cm}
    \begin{center}Table de vérité du \verb+or+\end{center}
    \begin{longtable}[]{@{}cc@{}}
        \hline Expression & Résultat\tabularnewline\hline\endhead{}
        True \texttt{or} True & True\tabularnewline{}
        True \texttt{or} False & True\tabularnewline{}
        False \texttt{or} True & True\tabularnewline{}
        False \texttt{or} False & False\tabularnewline\hline{}
    \end{longtable}
\end{minipage}

\pagebreak

\begin{easybox}{Booléen comme résultat de comparaison}{}
    Les booléens sont le résultat d'une comparaison. Les opérateurs de comparaison sont:
    \smallskip
    \begin{itemize}
        \begin{minipage}{0.44\linewidth}
            \item[] \texttt{==} \quad pour le test d'égalité
            \item[] \texttt{!=} \quad pour le test de différence
            \item[] \texttt{<} \quad pour le test d'infériorité
        \end{minipage}
        \begin{minipage}{0.44\linewidth}
            \item[] \texttt{>} \quad pour le test de supériorité
            \item[] \texttt{<=} \quad pour le test d'infériorité ou d'égalité
            \item[] \texttt{>=} \quad pour le test de supériorité ou d'égalité
        \end{minipage}
    \end{itemize}

    \end{easybox}


\begin{easybox}{Booléen comme résultat de comparaison}{}  
    \smallskip
    \begin{minipage}{0.32\linewidth}    
        \begin{ConsolePythontex}[ Label=false]{}
a = 5
b = 6
c = (a == b)
c
\end{ConsolePythontex}
    \end{minipage}\hfill
    \begin{minipage}{0.32\linewidth}
        \begin{ConsolePythontex}[ Label=false]{}
a = 5
b = 6
c = (a != b)
c
        \end{ConsolePythontex}
    \end{minipage}\hfill
    \begin{minipage}{0.32\linewidth}
        \begin{ConsolePythontex}[ Label=false]{}
a = 5
b = 6
c = (a < b)
c
        \end{ConsolePythontex}
    \end{minipage}
    
\end{easybox}


\begin{exercice_cours}{    Quelle est la valeur de la variable \verb+c+ après l'exécution des codes suivants?}{}


    \begin{minipage}[t]{0.4\linewidth}
\begin{lstlisting}[numbers=none]   
>>> a = 3
>>> b = 1
>>> c = a > b
\end{lstlisting}

\begin{lstlisting}[numbers=none]
>>> a = True
>>> b = False
>>> c = (a or b) and (not a)
\end{lstlisting}

    \end{minipage} \hfill
    \begin{minipage}[t]{0.58\linewidth}
\begin{lstlisting}[numbers=none]
>>> a, b = True, False
>>> c = (a and b) or (not a) 
\end{lstlisting}   

\begin{lstlisting}[numbers=none]
>>> a = True
>>> b = False
>>> x, y = 7, 9
>>> c = (b or (x == y)) or (a and (x != y))
\end{lstlisting}  
    \end{minipage}

\end{exercice_cours}
